\documentclass[11pt,a4paper]{article}
\usepackage[utf8]{inputenc}
\usepackage[T2A]{fontenc}
\usepackage[russian,english]{babel}
\usepackage{amsmath,amssymb,booktabs,array,longtable,geometry,xcolor,url}
\usepackage{tikz}
\usetikzlibrary{arrows.meta,calc,positioning,fit,shapes.geometric}
\geometry{top=1.6cm,bottom=1.6cm,left=1.8cm,right=1.5cm}
\setlength{\parindent}{0pt}
\setlength{\parskip}{5pt}
\newcommand{\code}[1]{\path{#1}}
\newcommand{\vect}[1]{\boldsymbol{#1}}
\title{Математическая постановка и фактическое исполнение GPU path tracer}
\author{Ревизия исходного кода проекта \code{path_trace_golang}}
\date{22 октября 2026 г.}
\begin{document}
\maketitle
\tableofcontents
\section{Назначение документа и границы точности}

Этот документ описывает не абстрактный «идеальный» path tracer, а текущую цепочку
исполнения проекта: выбор backend в \code{internal/usecases/render/service.go},
упаковку сцены в \code{internal/transport/gpu/packing.go}, compute shader из
\code{shaders.go}, вспомогательные shader-модули и обратную передачу результата.
Формулы ниже соответствуют реализации на 22.10.2026; приоритет имеют имена функций
и фактические операции шейдера.

CPU backend реализует ту же семантику трассировки в \code{float64}, но имеет
собственный worker pool и собственный RNG. Поэтому одинаковый seed делает каждый
backend воспроизводимым сам с собой, но не обязан давать попиксельно одинаковую
последовательность случайных чисел.

\subsection{Что было исправлено относительно старой версии}

\begin{longtable}{p{0.35\linewidth}p{0.56\linewidth}}
\toprule
Старое утверждение & Фактическое исполнение \\
\midrule
Work group $16\times16$ & Обычный path shader: $32\times32$; wireframe: $16\times16$. \\
RGBA16F & Текстуры и readback: RGBA32F; накопление: три float32 на пиксель. \\
Стратификация $4\times4$ & В path shader используется независимое jitter-сэмплирование; батч из 4 sample/pixel не является стратификацией. \\
Schlick Fresnel & Для диэлектрика используется точная неполяризованная формула Френеля; для металла --- Schlick-подобная цветная аппроксимация. \\
NEE только через $\cos\theta_i\cos\theta_o/(r^2p_A)$ & Сфера сэмплируется по телесному углу, коробка --- по видимой площади с преобразованием PDF. \\
Firefly clamp $L_{\max}=500$ & В текущем GPU path shader такого clamp нет. \\
Temporal accumulation через предыдущую текстуру & Текстура действительно ping-pong-ится, но shader не читает \code{prevFrameTex}; математическое накопление выполняет только \code{accumSSBO}. \\
\bottomrule
\end{longtable}

\section{Общая постановка задачи}

Пусть сцена задаётся множеством объектов $\mathcal O$, материалов $\mathcal M$,
источников света $\mathcal L$, камерой и средой. Для изображения ширины $W$ и
высоты $H$ требуется оценить цвет каждого пикселя $(x,y)$:
\[
 \widehat{C}(x,y)=\frac{1}{S}\sum_{k=0}^{S-1}L(x,y,k),
 \qquad S=\max(\texttt{SamplesPerPx},1).
\]
Здесь $L$ --- HDR-линейный RGB-радианс одного пути. Глубина пути ограничена
\code{MaxDepth}=D, а расстояние пересечения --- \code{MaxRayDist}; если оно не
положительно, GPU использует $R_{\max}=1000$.

Поверхностная часть оценивает уравнение рендеринга:
\[
 L_o(\vect{x},\omega_o)=L_e(\vect{x},\omega_o)+
 \int_{\Omega^+} f(\vect{x},\omega_i,\omega_o)L_i(\vect{x},\omega_i)
 (\vect{n}\cdot\omega_i)\,d\omega_i.
\]
В объёме добавляется затухание и single-scattering через коэффициенты
$\sigma_s$, $\sigma_a$, $\sigma_t=\sigma_s+\sigma_a$.

\section{Фактический pipeline: от вызова Go до PNG}

\begin{figure}[ht]
\centering
\begin{tikzpicture}[node distance=7mm and 7mm,>=Latex,
  block/.style={draw,rounded corners,fill=blue!7,align=center,minimum width=31mm,minimum height=9mm},
  data/.style={draw,rounded corners,fill=orange!12,align=center,minimum width=31mm,minimum height=9mm},
  arrow/.style={->,thick}]
\node[block] (service) {Service\\BackendGPU};
\node[block,right=of service] (worker) {GL worker\\locked OS thread};
\node[data,right=of worker] (scene) {Scene packing\\UBO/SSBO};
\node[block,below=of scene] (dispatch) {tiles +\\DispatchCompute};
\node[block,left=of dispatch] (shader) {compute shader\\one pixel/thread};
\node[data,left=of shader] (accum) {accumSSBO\\RGB float32};
\node[block,below=of accum] (readback) {PBO/readback\\HDR float};
\node[data,right=of readback] (post) {ACES + $\sqrt{\cdot}$\\dither/denoise};
\node[data,right=of post] (png) {image.RGBA\\PNG/UI};
\draw[arrow] (service)--(worker);
\draw[arrow] (worker)--(scene);
\draw[arrow] (scene)--(dispatch);
\draw[arrow] (dispatch)--(shader);
\draw[arrow] (shader)--(accum);
\draw[arrow] (accum)--(readback);
\draw[arrow] (readback)--(post);
\draw[arrow] (post)--(png);
\end{tikzpicture}
\caption{Реальная последовательность GPU backend.}
\end{figure}

\begin{enumerate}
\item \code{Service.IntoWithCancel} выбирает GPU только при \code{BackendGPU}; при
  ошибке инициализации/рендера выполняется CPU fallback.
\item \code{RenderWithCancel} создаёт скрытое GLFW-окно, запускает один GL worker и
  привязывает контекст к его OS thread. Нужен OpenGL 4.3 и compute shader.
\item \code{prepareGPUFrame} упаковывает сцену, изменяет размеры ресурсов, очищает
  texture и накопительный SSBO, загружает uniform/SSBO и выбирает normal/wireframe.
\item Один проход обрабатывает не весь кадр одним вызовом: изображение режется на
  адаптивные плитки. Каждая плитка ждёт fence, чтобы не переполнять очередь GPU и
  разрешать UI обрабатывать отмену.
\item После dispatch выполняется barrier, затем при live preview делается
  периодический readback через один из двух PBO. После последнего pass гарантируется
  финальный readback.
\item HDR RGB преобразуется в 8-bit RGBA, затем при подходящих настройках идут
  CPU bilateral denoise и необязательный box smooth.
\end{enumerate}

\section{Данные и раскладка памяти}

\subsection{Камера, небо и туман}

Все три UBO передаются как массивы по 16 float32, то есть по 64 байта:
\[
 C=[p_x,p_y,p_z,0,\ t_x,t_y,t_z,0,\ u_x,u_y,u_z,0,\mathrm{FOV},a,f,A],
\]
где $a$ --- aperture, $f$ --- focus distance, $A$ --- aspect ratio.

\[
 Sky=[s,0,0,0,\ c_r,c_g,c_b,1,\ h_r,h_g,h_b,1,\ z_r,z_g,z_b,1].
\]
Значение $s=0$ означает deprecated background, $s=1$ --- solid, $s=2$ ---
gradient. Для fog:
\[
 F=[d,q,a_s,v,\ c_r,c_g,c_b,1,\sigma_s,\sigma_a,g,k,n,o,0,0],
\]
где $d$ --- legacy density, $q$ --- scatter, $a_s$ --- affect sky, $v$ ---
GPU volumetric, $k$ --- heterogeneity strength, $n$ --- noise scale, $o$ ---
number of octaves.

\subsection{Материалы и объекты}

Для материала $i$ используется ровно 20 float32 (80 байт):
\[
 [type,rough,ior,smoothness,\ albedo_{rgb},reflectivity,\ emit_{rgb},pad,
 absorption_{rgb},absorptionScale,tint_{rgb},pad].
\]
При упаковке emission уже умножается на $\max(0,\texttt{Power})$. Для dielectric
при нулевых legacy-полях задаются IOR $1.5$, tint $(1,1,1)$ и absorption scale
$0.01$; для metal/mirror при нулевых полях smoothness и reflectivity становятся
соответственно $1-rough$ и $1$.

Для объекта используется 12 float32 (48 байт):
\[
 [type,materialIndex,0,0,\ position_{xyz},0,\ size_{xyz},0].
\]
Sphere использует $size_x$ как радиус, plane имеет нормаль $(0,1,0)$, box задаётся
\[
 \vect b_{min}=\vect p-\frac12\vect s,\qquad
 \vect b_{max}=\vect p+\frac12\vect s.
\]
Индекс источника хранится отдельно как int32 в \code{lightIndices}; туда попадают
только emissive-материалы с хотя бы одной положительной компонентой emission.

Накопитель имеет $3WH$ float32:
\[
 \mathrm{base}=3(yW+x),\qquad
 A_{\mathrm{base}:\mathrm{base}+3}\leftarrow A+L.
\]
RGBA32F texture и PBO имеют четыре float32 на пиксель, но четвёртая компонента
служит только альфой при записи $(R,G,B,1)$; в математическое накопление она не
входит.

\section{Параллелизм и индексы}

Для обычного shader объявлено $l_x=l_y=32$, поэтому одна work group содержит
$1024$ invocations. Для wireframe используется $16\times16$. Для участка плитки
размером $w\times h$:
\[
 G_x=\left\lceil\frac{w}{l_x}\right\rceil=\frac{w+l_x-1}{l_x},\qquad
 G_y=\left\lceil\frac{h}{l_y}\right\rceil.
\]
Глобальная координата потока с offset $(o_x,o_y)$:
\[
 x=\texttt{gl\_GlobalInvocationID.x}+o_x,\qquad
 y=\texttt{gl\_GlobalInvocationID.y}+o_y.
\]
Потоки за границей кадра сразу выходят по условию $x\ge W\lor y\ge H$.

Обычный path pass обрабатывает $B$ сэмплов на пиксель, где по умолчанию $B=4$ и
допустимы $1\ldots16$. Число pass:
\[
 P=\left\lceil\frac{S}{B}\right\rceil,
 \qquad B_i=\min(B,S-iB),
 \qquad N_i=iB+B_i.
\]
В shader каждый из $B_i$ sample использует индекс $k=N_i-B_i+j$. Начальный RNG:
\[
 q=yW+x,\qquad state_0=\operatorname{hash}(q\mathbin{\mathsf{xor}}
 \operatorname{hash}(frameSeed+k)).
\]
Затем каждый вызов \texttt{rng} делает
\[
 state\leftarrow hash(state+\texttt{0x9e3779b9}),\qquad
 \xi=\frac{state\mathbin{\gg}8}{2^{24}}\in[0,1).
\]
Здесь hash --- последовательность XOR, сдвигов и умножений на
\texttt{0xed5ad4bb}, \texttt{0xac4c1b51}, \texttt{0x31848bab}.

\section{Первичный луч камеры}

Пусть $\vect c$ --- position, $\vect t$ --- target, $\vect u_0$ --- up:
\[
 \vect w=\frac{\vect c-\vect t}{\|\vect c-\vect t\|},\qquad
 \vect u=\frac{\vect u_0\times\vect w}{\|\vect u_0\times\vect w\|},\qquad
 \vect v=\vect w\times\vect u.
\]
Для $\theta=\mathrm{FOV}\,\pi/180$:
\[
 h=\tan(\theta/2),\quad V_h=2h,\quad V_w=A V_h,
 \quad H=fV_h\vect v,\quad W_c=fV_w\vect u,
\]
\[
 \vect q=\vect c-\tfrac12W_c-\tfrac12H-f\vect w.
\]
Jitter выбирается без стратификации:
\[
 \xi_x,\xi_y\sim U[0,1),\qquad
 U=\frac{x+\xi_x}{W},\qquad
 V=\frac{H-1-y+\xi_y}{H}.
\]
Направление pinhole-луча:
\[
 \vect d=\operatorname{normalize}(\vect q+U W_c+V H-\vect c).
\]
При aperture $>0$ выбирается равномерная точка на диске объектива:
\[
 r_l=\frac{aperture}{2}\sqrt{\xi_1},\quad
 \phi=2\pi\xi_2,\quad
 \vect o=\vect c+\vect u(r_l\cos\phi)+\vect v(r_l\sin\phi),
\]
а направление строится из $\vect q+UW_c+VH-\vect o$. Это даёт DOF с фокусом
на плоскости, заданной $f$.

\begin{figure}[ht]
\centering
\begin{tikzpicture}[>=Latex,scale=1.05]
  \coordinate (c) at (0,0); \coordinate (p) at (5,0);
  \draw[thick,blue] (c) -- (p);
  \draw[thick,orange] (5,-1.5)--(5,1.5);
  \draw[->,very thick,purple] (c)--(4.1,0.7) node[midway,above] {$\vect o+t\vect d$};
  \draw[->,thick,red] (c)--(-1.1,0) node[left] {$\vect w$};
  \draw[->,thick,green!60!black] (c)--(0,1.1) node[right] {$\vect v$};
  \draw[->,thick,cyan!60!black] (c)--(0.8,0.7) node[above] {$\vect u$};
  \fill (c) circle (2pt) node[below left] {$\vect c$};
  \fill (p) circle (2pt) node[below right] {target/focus plane};
  \draw[dashed] (c)--(5,1.5); \draw[dashed] (c)--(5,-1.5);
\end{tikzpicture}
\caption{Построение луча из камеры; координата $y$ переворачивается перед вычислением $V$.}
\end{figure}

\section{Пересечения и геометрическая нормаль}

Все пересечения ищутся в интервале $[t_{min},t_{max}]$; для продолжения пути
\code{hitWorld} перебирает объекты в порядке массива и оставляет ближайшее
пересечение, каждый раз уменьшая $t_{max}$.

\subsection{Сфера}

Луч $\vect r(t)=\vect o+t\vect d$, сфера $(\vect c,r)$:
\[
 a=\vect d\cdot\vect d,\quad
 b_h=(\vect o-\vect c)\cdot\vect d,\quad
 c=(\vect o-\vect c)^2-r^2,
 \quad \Delta=b_h^2-ac.
\]
При $\Delta\ge0$ проверяются корни
\[
 t=\frac{-b_h\pm\sqrt{\Delta}}{a}.
\]
В текущем normal shader используется эквивалентная midpoint-форма:
\[
 t_m=-\frac{(\vect o-\vect c)\cdot\vect d}{a},\qquad
 \vect q=(\vect o-\vect c)+t_m\vect d,\qquad
 h=\sqrt{\frac{r^2-\vect q^2}{a}},\quad t=t_m-h.
\]
Если первый корень вне интервала, берётся $t_m+h$. Внешняя нормаль
$\vect n_o=\operatorname{normalize}(\vect r(t)-\vect c)$.

\subsection{Плоскость}

Для точки $\vect p$ и нормали $\vect n=(0,1,0)$:
\[
 t=\frac{(\vect p-\vect o)\cdot\vect n}{\vect d\cdot\vect n}.
\]
При нулевом знаменателе или $t\notin[t_{min},t_{max}]$ пересечения нет.

\subsection{AABB}

Для каждой оси $i$:
\[
 t_0^i=(b_{min}^i-o^i)/d^i,\quad
 t_1^i=(b_{max}^i-o^i)/d^i,
 \quad t_{near}=\max_i\min(t_0^i,t_1^i),
 \quad t_{far}=\min_i\max(t_0^i,t_1^i).
\]
Пересечение существует при $t_{far}\ge t_{near}$; если вход уже позади
$t_{min}$, берётся $t_{far}$. Для продолжения диэлектрического пути есть отдельный
режим \code{findExit}, который принудительно берёт выходную грань. Нормаль грани
получается из оси, обновившей $t_{near}$ или $t_{far}$, а не из сглаживания.

После пересечения применяется \code{setFaceNormal}:
\[
 frontFace=[\vect d\cdot\vect n_o<0],\qquad
 \vect n=frontFace?\vect n_o:-\vect n_o.
\]
Вторичный луч сдвигается от поверхности через ULP-based \code{offsetRay}; это
подавляет self-intersection и не является физическим расстоянием в сцене.

+
\section{Поверхности и throughput}

Для пути хранятся $\vect\beta$ (throughput), $\vect L$ (накопленная radiance),
текущая среда, $pdf_{prev}$, предыдущая точка и $\eta$-масштаб. Начальные значения:
\[
 \vect\beta=(1,1,1),\qquad \vect L=(0,0,0),\qquad etaScale=1.
\]

\subsection{Lambert}

\[
 f_r=\frac{\vect\rho}{\pi},\qquad
 p(\omega)=\frac{\vect n\cdot\omega}{\pi},\qquad
 \vect\beta\leftarrow\vect\beta\odot\vect\rho.
\]
Направление генерируется cosine-weighted sampling:
\[
 \phi=2\pi\xi_1,\quad \cos\theta=\sqrt{\xi_2},\quad
 \sin\theta=\sqrt{1-\xi_2},
\]
после чего локальный вектор преобразуется базисом $(\vect u,\vect v,\vect n)$.

\subsection{Metal: GGX visible-normal sampling}

Для metal shader использует $rough=1-smoothness$, $\alpha=rough^2$:
\[
 D(h)=\frac{\alpha^2}{\pi((1-n_h^2)+n_h^2\alpha^2)^2},
\]
\[
 G_1(c,\alpha)=\frac{2c}{c+\sqrt{\alpha^2+(1-\alpha^2)c^2}},
 \quad G=G_1(n_v,\alpha)G_1(n_l,\alpha),
\]
\[
 F=F_0+(1-F_0)(1-v_h)^5,
 \quad f=F\frac{DG}{4n_vn_l},
 \quad p_h=\frac{DG_1(n_v,\alpha)}{4n_v}.
\]
Вес нового направления:
\[
 \vect\beta\leftarrow\vect\beta\odot f\,\frac{n_l}{p_h}.
\]
Сэмплирование выбирает видимую микронормаль через stretched view vector и
disk-параметризацию; это не простое равномерное сэмплирование GGX.

\subsection{Mirror и гладкий metal}

\[
 \vect\omega_r=\vect d-2(\vect d\cdot\vect n)\vect n,
 \qquad \vect\beta\leftarrow\vect\beta\odot\vect{albedo}\,reflectivity.
\]
Для metal дополнительно применяется цветной $F$ из предыдущей формулы. Mirror
не получает direct-light sample и отражает идеально.

\subsection{Dielectric: интерфейс сред}

Текущая среда определяется не только front/back флагом: shader ищет ближайший
dielectric object, содержащий точку после ULP-offset. Поэтому вложенные стеклянные
объекты имеют стекоподобную семантику.

\[
 \eta=\frac{\eta_1}{\eta_2},\quad
 c=\max(-\vect d\cdot\vect n,0),\quad
 s_t^2=\eta^2(1-c^2).
\]
При $s_t^2>1$ возникает TIR и используется reflection. Иначе:
\[
 \vect r_\perp=\eta(\vect d+c\vect n),\qquad
 \vect r_\parallel=-\sqrt{1-\|\vect r_\perp\|^2}\vect n,
 \quad \vect d_t=\vect r_\perp+\vect r_\parallel.
\]
Вероятность отражения --- точные уравнения Френеля:
\[
 s_t^2=\eta^2(1-c^2),\quad c_t=\sqrt{1-s_t^2},
\]
\[
 R=\frac12\left[
 \left(\frac{\eta c-c_t}{\eta c+c_t}\right)^2+
 \left(\frac{c-\eta c_t}{c+\eta c_t}\right)^2\right].
\]
При преломлении $\vect\beta\leftarrow\eta^2\vect\beta$ и обновляется
$etaScale\leftarrow etaScale/\eta^2$; при отражении вес равен $(1,1,1)$.
На пути внутри dielectric прозрачность дополнительно умножается на
\[
 T_{glass}(d)=\exp(-\vect\sigma\,scale\,d)\odot
 \vect{tint}^{\,d}.
\]

\section{Прямое освещение и MIS}

\subsection{Список света}

Источники выбираются из emissive sphere и box. За одну вершину shader samples
$K=\min(|\mathcal L|,8)$ источников; если источников больше восьми, стартовый индекс
случайный и берётся циклическое окно. Inclusion factor:
\[
 q=\min\left(1,\frac{8}{|\mathcal L|}\right).
\]

Для spherical light точка выбирается по видимому телесному углу. В терминах центра
на расстоянии $d$ и радиуса $r$:
\[
 s^2=\frac{r^2}{d^2},\quad 1-\cos\theta=\frac{s^2}{1+\sqrt{1-s^2}},
 \quad p_\Omega=\frac{1}{2\pi(1-\cos\theta)}.
\]
Для box видимые грани имеют площади
\[
 A_x=s_ys_z,\quad A_y=s_xs_z,\quad A_z=s_xs_y;
\]
грань обнуляется, если reference находится внутри её slab. После выбора точки
равномерно по суммарной площади $A$:
\[
 p_\Omega=\frac{r^2}{A\,|\vect n_L\cdot(-\vect\omega)|}.
\]

Для shadow ray проверяется occlusion всеми объектами, кроме самого emitter; затем
учитывается glass/fog transmittance. Поверхностный direct contribution:
\[
 \Delta L=\vect f\odot\vect L_e\odot T\;
 \frac{n_l}{p_Lq+[\text{continuation}]p_B}.
\]
При первом/последнем bounce $p_B$ не добавляется, иначе denominator реализует
MIS balance heuristic между light sampling и BSDF sampling.

При прямом попадании случайного пути в источник его вклад умножается на
\[
 w_e=\frac{p_{prev}}{p_{prev}+p_Lq},
\]
а если предыдущий PDF нулевой, $w_e=1$.

\begin{figure}[ht]
\centering
\begin{tikzpicture}[>=Latex,scale=1.0]
  \draw[thick,blue] (-3,0)--(3,0) node[right] {surface};
  \fill (0,0) circle (2pt) node[below left] {$x$};
  \draw[->,thick,green!60!black] (0,0)--(0,1.4) node[right] {$n$};
  \draw[thick,orange,fill=orange!25] (1.8,2.1) circle (0.45);
  \fill (1.8,2.1) circle (2pt) node[above] {light};
  \draw[->,thick,red] (0,0)--(1.8,2.1) node[midway,left] {$\omega_i$};
  \draw[dashed] (0,0)--(1.8,1.65);
  \node at (0,-0.65) {$\Delta L\propto f\,L_e\,T\,n_l/(p_Lq+p_B)$};
\end{tikzpicture}
\caption{Next Event Estimation: отдельный shadow ray к сэмплу источника.}
\end{figure}

\section{Итеративный path loop}

Для каждого bounce $b=0,\ldots,D-1$ shader выполняет следующие действия:

\begin{enumerate}
\item Проверяет $\max(\beta_r,\beta_g,\beta_b)>0$.
\item Ищет ближайшее пересечение до $R_{\max}$ и получает расстояние $d$.
\item Если луч ещё не внутри стекла и включён volumetric fog, sample-ит возможное
  рассеяние до ближайшей поверхности/неба. При событии изменяет throughput, добавляет
  direct volume light, сэмплирует фазу и продолжает путь из точки объёма.
\item Иначе умножает throughput на стеклянное затухание текущей среды.
\item При miss добавляет $\vect\beta\odot background(\vect d)$ и завершает путь.
  Для legacy non-volumetric fog с affect-sky фон дополнительно проходит
  $F(L,d)=Le^{-density\,d}+C(1-e^{-density\,d})$ при $d=50$.
\item При emissive hit добавляет $\vect\beta\odot emit\cdot w_e$ только для
  front face и завершает путь.
\item Для Lambert/metal/mirror/dielectric вычисляет direct/BSDF/refraction,
  сохраняет $p_{prev}$ и previous point, умножает throughput и строит offset ray.
\item После $b\ge3$ запускает Russian roulette:
  \[
   p_{rr}=\min(1,etaScale\max(\beta_r,\beta_g,\beta_b)).
  \]
  При отказе путь завершается, при продолжении $\vect\beta\leftarrow\vect\beta/p_{rr}$.
\end{enumerate}

\section{Объёмный туман}

Если явные коэффициенты не заданы, из legacy density $d$ и scatter $q$ строится:
\[
 \sigma_s=dq,\qquad \sigma_a=\max(d-\sigma_s,0),\qquad
 \sigma_t=\sigma_s+\sigma_a.
\]
Для heterogeneous fog коэффициенты умножаются на
\[
 m(\vect x)=(1-k)+2kN(\vect x),
\]
где $N$ --- hash/fBm noise с 1--5 октавами.

GPU выбирает majorant:
\[
 \mu=\bigl(\max(\sigma_s,0)+\max(\sigma_a,0)\bigr)(1+k),
\]
с fallback на density, если физические коэффициенты нулевые. Кандидат дистанции
delta tracking:
\[
 t\leftarrow t-\frac{\ln(1-\xi)}{\mu}.
\]
Если $t\ge d$, событие объёма не произошло. Иначе принимается с вероятностью
$\sigma_t(\vect x)/\mu$. Вес принятого scattering event:
\[
 \vect w_{fog}=\operatorname{clamp}(fogColor,0,1)\frac{\sigma_s(\vect x)}{\sigma_t(\vect x)}.
\]
Для неоднородной shadow transmittance применяется ratio tracking:
\[
 T\leftarrow T\left(1-\frac{\sigma_t(\vect x)}{\mu}\right)
\]
на каждой случайной позиции до конца отрезка. В однородном случае shader использует
аналитическое $T=\exp(-\mu d)$.

Фаза Хеньи--Грина:
\[
 p_{HG}(\cos\theta)=\frac{1-g^2}{4\pi(1+g^2-2g\cos\theta)^{3/2}}.
\]
Volume NEE использует ту же схему максимум восьми источников, а contribution имеет
вид
\[
 \Delta L_{vol}=L_e\frac{p_{HG}}{p_Lq+[\text{continuation}]p_{HG}}T_{shadow}.
\]

\begin{figure}[ht]
\centering
\begin{tikzpicture}[>=Latex,scale=1.0]
  \draw[thick,blue,fill=blue!8] (-2,-1)--(3,-1)--(3,1.4)--(-2,1.4)--cycle;
  \node[blue] at (0.5,1.05) {heterogeneous medium};
  \fill (-2.3,0) circle (2pt) node[left] {camera};
  \draw[->,very thick,red] (-2.3,0)--(3.3,0);
  \foreach \x in {-1.5,-0.7,0.1,0.9,1.7,2.5}{
    \fill[purple] (\x,0) circle (2pt);
    \draw[->,thin,purple,dashed] (\x,0)--(2.0,1.0);
  }
  \fill[orange] (2.0,1.0) circle (5pt) node[above right] {emitter};
  \node at (0,-1.45) {$t\sim\mathrm{Exp}(\mu)$, accept with $\sigma_t/\mu$, weight $\sigma_s/\sigma_t$};
\end{tikzpicture}
\caption{Delta tracking вдоль луча и прямое освещение точки рассеяния.}
\end{figure}

+
\section{Накопление, HDR и постобработка}

После $B_i$ путей одного pass для пикселя:
\[
 A_{x,y}\leftarrow A_{x,y}+\sum_{j=0}^{B_i-1}L_{x,y,k+j},
 \qquad \overline L_{x,y}=\frac{A_{x,y}}{N_i}.
\]
Именно $N_i$ передаётся как \code{uSampleCount}; поэтому последний неполный batch
также усредняется корректно.

\subsection{Тонемаппинг и 8-bit output}

Для каждой положительной HDR-компоненты:
\[
 A(x)=\operatorname{clamp}\left(
 \frac{x(2.51x+0.03)}{x(2.43x+0.59)+0.14},0,1\right),
 \qquad C_{8}=255\sqrt{A(x)}+dither.
\]
NaN, infinity и отрицательные значения становятся нулём. Альфа всегда равна 255.
Это ACES-like mapping плюс square-root gamma; это не linear-to-sRGB transfer function.

\subsection{CPU denoise и smooth}

При включённом denoise и $S\ge4$ после readback запускаются три прохода bilateral
фильтра на радиусах $1,2,4$. Для соседнего пикселя $j$:
\[
 w_{sp}=\exp\left(-\frac{\Delta x^2+\Delta y^2}{2\sigma_s^2}\right),\quad
 b=\frac{1}{2\max(\sigma_r^2\,255^2,10^{-4})},
\]
\[
 w=\frac{w_{sp}}{\left(1+(\Delta R^2+\Delta G^2+\Delta B^2)b\right)^2},
 \qquad C'=\frac{\sum_jwC_j}{\sum_jw}.
\]
Дополнительный smooth выключен по умолчанию; если включён, его результат смешивается
с исходным byte-color с strength и радиусом 1--5.

\section{Фактические настройки исполнения}

\begin{longtable}{p{0.38\linewidth}p{0.52\linewidth}}
\toprule
Параметр & Реальное значение/правило \\
\midrule
Normal work group & $32\times32$ \\
Wireframe work group & $16\times16$ \\
Samples per pass & default $4$, clamp $1\ldots16$ \\
Dispatch tile & старт 64, предел по умолчанию 256, диапазон 32--512; адаптируется по времени fence \\
Max light samples & $8$ \\
Max smooth radius & $5$ \\
GPU texture & RGBA32F \\
Accumulation & SSBO: $3WH$ float32 \\
Readback & два PBO, RGBA float32 \\
Temporal flag & target texture ping-pong есть, но shader не вызывает texture lookup предыдущего кадра \\
Denoise & CPU, по умолчанию включён для GPU при $S\ge4$ \\
GPU fallback & при ошибке GPU service запускает CPU renderer \\
\bottomrule
\end{longtable}

\section{Проверка соответствия коду}

Ниже перечислены точки, которые следует проверять при изменении математики:

\begin{longtable}{p{0.37\linewidth}p{0.53\linewidth}}
\toprule
Математический блок & Источник истины в проекте \\
\midrule
Backend/fallback & \code{internal/usecases/render/service.go}: \code{IntoWithCancel}, \code{renderIntoGPU} \\
Packing/stride & \code{internal/transport/gpu/packing.go}: \code{packMaterials}, \code{packObjects}, \code{packFogBlock} \\
GL resources/readback & \code{internal/transport/gpu/setup.go}, \code{gpu.go}: \code{resizeGPUResources}, \code{tonemapHDRToRGBA} \\
Dispatch/passes & \code{internal/transport/gpu/samples.go}: \code{samplePasses}, \code{dispatchTiles}, \code{dispatchGroups} \\
Camera/RNG/path & \code{internal/transport/gpu/shaders.go}, \code{physics_shader.go}, \code{transport_shader.go} \\
Lights/BSDF/MIS & \code{lighting_shader.go}, \code{physics_shader.go} \\
Fog & \code{fog_shader.go}, \code{physics_shader.go} \\
CPU parity & \code{internal/usecases/render/cpu_gpu_parity.go}, \code{physics.go}, \code{fog_physics.go} \\
Validation & \code{scenes/validation/README.md} и GPU internal tests \\
\bottomrule
\end{longtable}

Для регрессионной проверки применяются команды:
\begin{verbatim}
go test ./internal/transport/gpu ./internal/usecases/render ./internal/service
go run ./cmd/validate-render -scene scenes/validation/glass_grid.json \
  -backend gpu -width 320 -height 180 -spp 64 -seed 1 -out bin/glass.png
PATHTRACER_GPU_INTEGRATION=1 go test ./internal/transport/gpu \
  -run '^TestGPU(PhysicsAudit|Integration)$' -count=1 -v -timeout=180s
\end{verbatim}

В среде без аппаратного GL последний контур может быть недоступен; намеренное
использование software Mesa требует \code{PATHTRACER\_GPU\_ALLOW\_SOFTWARE=1}.
MSE/PSNR после ACES и gamma измеряют качество изображения, но не доказывают
точность линейного интеграла.

\section{Итоговая математическая схема}

Для пикселя и sample $k$ фактическая оценка имеет вид:
\[
 L_k=\operatorname{PathTrace}\left(
 \operatorname{CameraRay}(x,y,\xi_k),
 \mathcal O,\mathcal M,\mathcal L,\mathcal F,D,R_{\max}
 \right),
\]
где внутри каждого bounce:
\[
 \begin{aligned}
 &\text{hit/medium event}\\
 &\beta\leftarrow\beta\odot T_{glass/fog}\\
 &L\leftarrow L+\beta\odot(L_e+L_{direct})\\
 &\omega_{next}\sim p_{Lambert/GGX/Fresnel/HG}\\
 &\beta\leftarrow\beta\odot\frac{f\,|n\cdot\omega_{next}|}{p}
   \quad\text{(где применимо)}\\
 &\text{Russian roulette после третьего bounce}.
 \end{aligned}
\]
Затем GPU выполняет:
\[
 A\leftarrow A+\sum_kL_k,\qquad
 \overline L=\frac{A}{N},\qquad
 C_{8}=255\sqrt{\operatorname{ACES}(\overline L)}+dither.
\]
Именно последняя формула, а не clamp/Firefly reduction, является текущей границей
между линейным HDR накопителем и отображаемым изображением.

\end{document}

\end{document}
